% Pembaca lengkap Open Logic Project dalam Bahasa Indonesia.
% Sumber Inggris beku: OpenLogicProject/OpenLogic
% 9620cc73f9c8e0ad003c514a5d3748f29611c4c0.

\documentclass[openany]{memoir}

% `cleveref` Open Logic belum memiliki opsi Indonesian. Babel tetap memuat
% Indonesian; `\ollanguage` bernilai english hanya selama cleveref dimuat.
\PassOptionsToPackage{english,indonesian}{babel}

\newcommand*{\olpath}{../..}
\newcommand*{\ollangid}{id}
\newcommand*{\ollanguage}{english}

\usepackage{mathpazo}
\usepackage{helvet}
\usepackage{courier}
\linespread{1.05}

\usepackage{gitinfo2}

\input{\olpath/sty/open-logic.sty}
\renewcommand*{\ollanguage}{indonesian}
\input{\olpath/sty/open-logic-defer.sty}

\hypersetup{
  pdftitle={Open Logic Project — Edisi Lengkap Bahasa Indonesia (OLP-0722)},
  pdfauthor={Open Logic Project; terjemahan oleh OpenAI 5.6 Sol, Ultra mode},
  pdfsubject={Terjemahan Bahasa Indonesia dari Open Logic Project pada revisi beku 9620cc73f9c8e0ad003c514a5d3748f29611c4c0},
  pdfkeywords={Open Logic Project, Bahasa Indonesia, logika matematika, terjemahan, OLP-0722}
}

% `gitinfo2` emits `Release: (None) ((None))` when the checkout has no
% generated Git metadata.  Bind both its memoir footer and OLP's short
% revision text deterministically to the frozen source instead.
\makeatletter
\renewcommand*{\@gitFootRev}{Revisi sumber: 9620cc7 (12-07-2026)}
\makeatother
\renewcommand*{\ifgitinfo}{\sffamily\footnotesize
  revisi~9620cc7~(12-07-2026)}

\includeenv{editorial}
\problemsperchapter
\let\cleardoublepage\clearpage
\pagestyle{plain}

\begin{document}

\selectlanguage{indonesian}

\begin{titlingpage}
\begin{raggedleft}
\HUGE
\selectfont\bfseries\sffamily
\MakeUppercase{\OLTname}
\vskip 3ex
\normalfont\huge
Edisi Lengkap Bahasa Indonesia
\vskip 4ex
\textbf{\href{http://openlogicproject.org/}{\OLPname}}
\vskip 4ex
\Large
Revisi sumber beku:\\
\texttt{9620cc73f9c8e0ad003c514a5d3748f29611c4c0}\\
12 Juli 2026
\vskip 2ex
\large
Versi \texttt{OLP-0722-20260814}\\
722/722 berkas isi diterjemahkan\\
\small
642 modul dalam pembaca; 80 modul nonpembaca dipertahankan dalam sumber
\end{raggedleft}

\vfill

\ollicense

\vskip 2ex
\footnotesize
Edisi ini adalah \textbf{terjemahan Bahasa Indonesia dari} (DataCite:
\emph{IsTranslationOf}) sumber Inggris Open Logic Project pada revisi beku
\href{https://github.com/OpenLogicProject/OpenLogic/commit/9620cc73f9c8e0ad003c514a5d3748f29611c4c0}%
{\texttt{9620cc73f9c8e0ad003c514a5d3748f29611c4c0}}.
DOI rilis: \href{https://doi.org/10.5281/zenodo.21932787}%
{\texttt{10.5281/zenodo.21932787}}; DOI konsep Bahasa Indonesia:
\href{https://doi.org/10.5281/zenodo.21932786}%
{\texttt{10.5281/zenodo.21932786}}. Repositori:
\href{https://github.com/KokunoYumeto/OpenLogic-id}%
{\texttt{github.com/KokunoYumeto/OpenLogic-id}}.
Seluruh perubahan bahasa dan koreksi sumber terbatas dicatat dalam sumber
pelokalan dan provenans yang menyertainya. Open Logic Project tidak menyokong
atau mengesahkan terjemahan ini.

\end{titlingpage}

\pagestyle{giheadings}

\chapter*{Tentang Proyek Logika Terbuka}
\addcontentsline{toc}{chapter}{Tentang Proyek Logika Terbuka}


\textit{Teks Logika Terbuka} adalah buku ajar sumber terbuka dan
kolaboratif tentang metalogika formal dan metode formal, yang dimulai pada
tingkat menengah (yakni, setelah mata kuliah pengantar logika formal).
Walaupun ditujukan kepada pembaca berlatar nonmatematika (khususnya mahasiswa
filsafat dan ilmu komputer), buku ini menyajikan pembahasan yang ketat.

Cakupan beberapa topik yang saat ini disertakan mungkin belum lengkap, dan
banyak bagian masih memerlukan revisi besar. Kami berencana memperluas teks
ini agar kelak mencakup lebih banyak topik. Kami juga berencana menambahkan
sejumlah unsur, seperti glosarium, daftar bacaan lanjutan, catatan historis,
gambar, penjelasan yang lebih baik, bagian yang menjelaskan relevansi
hasil-hasil bagi filsafat, ilmu komputer, dan matematika, serta lebih banyak
soal dan contoh. Jika Anda menemukan kesalahan atau memiliki saran,
\href{https://github.com/OpenLogicProject/OpenLogic/wiki/Contributing}{silakan
beri tahu tim proyek}.

Proyek ini berjalan dalam semangat sumber terbuka. Teks ini tidak hanya
tersedia secara bebas; kami juga menyediakan sumber LaTeX berdasarkan
Lisensi Creative Commons Atribusi, yang memberi setiap orang hak untuk
mengunduh, menggunakan, mengubah, menyusun ulang, mengonversi, dan
mendistribusikan kembali karya kami, asalkan mereka memberikan atribusi yang
sesuai. Informasi tambahan tersedia di situs web Proyek Logika Terbuka,
\href{http://openlogicproject.org/}{openlogicproject.org}.


\tableofcontents*

\olimport[content]{content}

\photocredits

\bibliographystyle{\olpath/bib/natbib-oup}
\bibliography{\olpath/bib/open-logic}

\end{document}
